Here is a self-contained corrected proof.

```latex
\[
\newcommand{\R}{\mathbb{R}}
\newcommand{\diag}{\operatorname{diag}}
\]

For \(r \geq 1\), write
\[
\|A\|_{r\to r}:=\sup_{x\neq 0}\frac{\|Ax\|_r}{\|x\|_r}
\]
for the induced \(\ell^r\)-operator norm, and define
\[
\rho_r(A):=\inf\left\{\|DAD^{-1}\|_{r\to r}: D \text{ is positive diagonal}\right\}.
\]
Here positive diagonal means diagonal with all diagonal entries strictly positive.

\[
\textbf{Part (a).}
\]

Let
\[
s:=n-m.
\]
We first prove the block norm identity
\[
\|\diag(A,0_s)\|_{r\to r}=\|A\|_{r\to r}
\]
for every \(A\in \R^{m\times m}\) and every \(s\geq 0\).

Take \(x=(x_1,x_2)\in \R^m\oplus \R^s\). Since \(r\geq 1\),
\[
\|x\|_r^r=\|x_1\|_r^r+\|x_2\|_r^r.
\]
Hence
\[
\|x_1\|_r^r\leq \|x\|_r^r,
\]
and therefore
\[
\|x_1\|_r\leq \|x\|_r.
\]
Thus
\[
\|\diag(A,0_s)x\|_r
=
\|(Ax_1,0)\|_r
=
\|Ax_1\|_r
\leq
\|A\|_{r\to r}\|x_1\|_r
\leq
\|A\|_{r\to r}\|x\|_r.
\]
Taking the supremum over \(x\neq 0\) gives
\[
\|\diag(A,0_s)\|_{r\to r}\leq \|A\|_{r\to r}.
\]

For the reverse inequality, the unit sphere
\[
\{u\in \R^m:\|u\|_r=1\}
\]
is compact, and the function \(u\mapsto \|Au\|_r\) is continuous. Therefore there exists \(u^\ast\in \R^m\) with \(\|u^\ast\|_r=1\) such that
\[
\|Au^\ast\|_r=\|A\|_{r\to r}.
\]
Testing \(\diag(A,0_s)\) on \((u^\ast,0)\) gives
\[
\|\diag(A,0_s)\|_{r\to r}
\geq
\frac{\|\diag(A,0_s)(u^\ast,0)\|_r}{\|(u^\ast,0)\|_r}
=
\frac{\|(Au^\ast,0)\|_r}{\|u^\ast\|_r}
=
\|Au^\ast\|_r
=
\|A\|_{r\to r}.
\]
Hence
\[
\|\diag(A,0_s)\|_{r\to r}=\|A\|_{r\to r}.
\]

Now let \(D\in \R^{n\times n}\) be positive diagonal. With respect to the decomposition
\[
\R^n=\R^m\oplus \R^{n-m},
\]
we may write
\[
D=\diag(D_1,D_2),
\]
where \(D_1\in \R^{m\times m}\) and \(D_2\in \R^{(n-m)\times (n-m)}\) are positive diagonal. If \(n=m\), the \(D_2\)-block is simply absent.

Then
\[
D M_n D^{-1}
=
\diag(D_1,D_2)
\begin{pmatrix}
H&0\\
0&0_{n-m}
\end{pmatrix}
\diag(D_1^{-1},D_2^{-1})
=
\diag(D_1HD_1^{-1},0_{n-m}).
\]
By the block norm identity,
\[
\|D M_n D^{-1}\|_{r\to r}
=
\|D_1HD_1^{-1}\|_{r\to r}.
\]
As \(D\) ranges over all positive diagonal \(n\times n\) matrices, \(D_1\) ranges over all positive diagonal \(m\times m\) matrices, while the expression above is independent of \(D_2\). Therefore
\[
\rho_r(M_n)
=
\inf_{D}\|D M_n D^{-1}\|_{r\to r}
=
\inf_{D_1}\|D_1HD_1^{-1}\|_{r\to r}
=
\rho_r(H).
\]
Thus
\[
\boxed{\rho_r(M_n)=\rho_r(H)}
\]
for every \(r\geq 1\).

\[
\textbf{Part (b).}
\]

We first prove the Hadamard norm formula needed below.

Let \(H\in \R^{m\times m}\) be a Hadamard matrix, meaning that all entries of \(H\) are \(\pm 1\) and
\[
HH^{\mathsf T}=mI_m.
\]
Since \(HH^{\mathsf T}=mI_m\), the matrix \(H\) is invertible, and
\[
H^{-1}=\frac{1}{m}H^{\mathsf T}.
\]
Multiplying by \(H\) on the right gives
\[
H^{\mathsf T}H=mI_m.
\]

We claim that for every \(q\geq 2\),
\[
\rho_q(H)=m^{1-\frac1q}.
\]

First we prove the upper bound. We use the standard finite-dimensional norm comparison: if \(1\leq a\leq q\), then for every \(z\in \R^m\),
\[
\|z\|_a\leq m^{\frac1a-\frac1q}\|z\|_q.
\]
Indeed, for \(a<q\), Hölder's inequality with conjugate exponents \(q/a\) and \(q/(q-a)\) gives
\[
\sum_{j=1}^m |z_j|^a
\leq
\left(\sum_{j=1}^m |z_j|^q\right)^{a/q}
\left(\sum_{j=1}^m 1\right)^{1-a/q}
=
\|z\|_q^a m^{1-a/q}.
\]
Taking \(a\)-th roots gives
\[
\|z\|_a\leq m^{\frac1a-\frac1q}\|z\|_q.
\]
The case \(a=q\) is equality.

Let
\[
C:=m^{1-\frac1q}.
\]
If \(q=2\), then for every \(x\in \R^m\),
\[
\|Hx\|_2^2
=
x^{\mathsf T}H^{\mathsf T}Hx
=
m\|x\|_2^2,
\]
so
\[
\|H\|_{2\to 2}=\sqrt m=m^{1-\frac12}.
\]

Now assume \(q>2\). For \(x\in \R^m\), set
\[
y:=Hx.
\]
Using \(H^{\mathsf T}H=mI_m\), we get
\[
\|y\|_2=\|Hx\|_2=\sqrt m\,\|x\|_2.
\]
By the norm comparison with \(a=2\),
\[
\|x\|_2\leq m^{\frac12-\frac1q}\|x\|_q.
\]
Hence
\[
\|y\|_2\leq \sqrt m\,m^{\frac12-\frac1q}\|x\|_q
=
m^{1-\frac1q}\|x\|_q
=
C\|x\|_q.
\]
Also, since every entry of \(H\) is \(\pm 1\),
\[
|(Hx)_i|
\leq
\sum_{j=1}^m |x_j|
=
\|x\|_1.
\]
Using the norm comparison with \(a=1\),
\[
\|x\|_1\leq m^{1-\frac1q}\|x\|_q
=
C\|x\|_q.
\]
Therefore
\[
\|y\|_\infty=\|Hx\|_\infty\leq C\|x\|_q.
\]

Because \(q>2\), there is no \(0^0\) issue in the following estimate. For each coordinate,
\[
|y_i|^q=|y_i|^{q-2}|y_i|^2
\leq
\|y\|_\infty^{q-2}|y_i|^2.
\]
Summing over \(i\) gives
\[
\|y\|_q^q
=
\sum_{i=1}^m |y_i|^q
\leq
\|y\|_\infty^{q-2}\sum_{i=1}^m |y_i|^2
=
\|y\|_\infty^{q-2}\|y\|_2^2.
\]
Using the two bounds above,
\[
\|Hx\|_q^q
=
\|y\|_q^q
\leq
(C\|x\|_q)^{q-2}(C\|x\|_q)^2
=
C^q\|x\|_q^q.
\]
Taking \(q\)-th roots,
\[
\|Hx\|_q\leq C\|x\|_q.
\]
Thus
\[
\|H\|_{q\to q}\leq m^{1-\frac1q}.
\]
Since \(I_m\) is an admissible positive diagonal matrix in the definition of \(\rho_q(H)\),
\[
\rho_q(H)\leq \|H\|_{q\to q}\leq m^{1-\frac1q}.
\]

Now we prove the reverse inequality. Let
\[
D=\diag(d_1,\dots,d_m)
\]
be any positive diagonal matrix. Choose \(i\) such that
\[
d_i=\max_{1\leq j\leq m} d_j.
\]
Let
\[
u:=H^{\mathsf T}e_i.
\]
Then \(u\in \{\pm 1\}^m\), because the entries of \(H\) are \(\pm 1\), and
\[
Hu=HH^{\mathsf T}e_i=m e_i.
\]
Set
\[
x:=Du.
\]
Then \(x\neq 0\), and
\[
DHD^{-1}x
=
DHD^{-1}Du
=
DHu
=
D(me_i)
=
m d_i e_i.
\]
Therefore
\[
\|DHD^{-1}x\|_q=m d_i.
\]
On the other hand,
\[
\|x\|_q
=
\|Du\|_q
=
\left(\sum_{j=1}^m d_j^q |u_j|^q\right)^{1/q}
=
\left(\sum_{j=1}^m d_j^q\right)^{1/q}
\leq
\left(\sum_{j=1}^m d_i^q\right)^{1/q}
=
m^{1/q}d_i.
\]
Hence
\[
\|DHD^{-1}\|_{q\to q}
\geq
\frac{\|DHD^{-1}x\|_q}{\|x\|_q}
\geq
\frac{m d_i}{m^{1/q}d_i}
=
m^{1-\frac1q}.
\]
Since this holds for every positive diagonal \(D\),
\[
\rho_q(H)\geq m^{1-\frac1q}.
\]
Combining the two inequalities,
\[
\boxed{\rho_q(H)=m^{1-\frac1q}}
\]
for every \(q\geq 2\).

It remains to exhibit suitable Hadamard matrices of orders \(4\) and \(8\).

Let
\[
H_4=
\begin{pmatrix}
1&1&1&1\\
1&-1&1&-1\\
1&1&-1&-1\\
1&-1&-1&1
\end{pmatrix}.
\]
Let its rows be \(r_1,r_2,r_3,r_4\). Each row has squared norm \(4\). The off-diagonal inner products are
\[
\begin{aligned}
r_1\cdot r_2&=1-1+1-1=0,\\
r_1\cdot r_3&=1+1-1-1=0,\\
r_1\cdot r_4&=1-1-1+1=0,\\
r_2\cdot r_3&=1-1-1+1=0,\\
r_2\cdot r_4&=1+1-1-1=0,\\
r_3\cdot r_4&=1-1+1-1=0.
\end{aligned}
\]
Thus
\[
H_4H_4^{\mathsf T}=4I_4,
\]
so \(H_4\) is Hadamard.

Now define
\[
H_8:=
\begin{pmatrix}
H_4&H_4\\
H_4&-H_4
\end{pmatrix}.
\]
All entries of \(H_8\) are \(\pm 1\). Moreover,
\[
H_8H_8^{\mathsf T}
=
\begin{pmatrix}
H_4&H_4\\
H_4&-H_4
\end{pmatrix}
\begin{pmatrix}
H_4^{\mathsf T}&H_4^{\mathsf T}\\
H_4^{\mathsf T}&-H_4^{\mathsf T}
\end{pmatrix}.
\]
Multiplying the block matrices,
\[
H_8H_8^{\mathsf T}
=
\begin{pmatrix}
2H_4H_4^{\mathsf T}&0\\
0&2H_4H_4^{\mathsf T}
\end{pmatrix}
=
\begin{pmatrix}
8I_4&0\\
0&8I_4
\end{pmatrix}
=
8I_8.
\]
Therefore \(H_8\) is Hadamard.

Now let \(n>6\) and let \(p>2\) be an integer. Then \(p\geq 3\), so both \(p\) and \(p+1\) are at least \(2\), and the Hadamard formula above applies.

\[
\textbf{Case 1: } n=7.
\]
Take
\[
M=\diag(H_4,0_3)\in \R^{7\times 7}.
\]
By part \((a)\),
\[
\rho_p(M)=\rho_p(H_4)
\qquad\text{and}\qquad
\rho_{p+1}(M)=\rho_{p+1}(H_4).
\]
Using the formula \(\rho_q(H_4)=4^{1-1/q}\), we get
\[
\rho_p(M)=4^{1-\frac1p}
\]
and
\[
\rho_{p+1}(M)=4^{1-\frac1{p+1}}.
\]
Since
\[
\frac1p>\frac1{p+1},
\]
we have
\[
1-\frac1p<1-\frac1{p+1}.
\]
Because \(4>1\), the function \(t\mapsto 4^t\) is strictly increasing. Therefore
\[
4^{1-\frac1p}<4^{1-\frac1{p+1}},
\]
so
\[
\rho_p(M)<\rho_{p+1}(M).
\]

\[
\textbf{Case 2: } n\geq 8.
\]
Take
\[
M=\diag(H_8,0_{n-8})\in \R^{n\times n},
\]
with the zero block omitted when \(n=8\). By part \((a)\),
\[
\rho_p(M)=\rho_p(H_8)
\qquad\text{and}\qquad
\rho_{p+1}(M)=\rho_{p+1}(H_8).
\]
Using the formula \(\rho_q(H_8)=8^{1-1/q}\), we obtain
\[
\rho_p(M)=8^{1-\frac1p}
\]
and
\[
\rho_{p+1}(M)=8^{1-\frac1{p+1}}.
\]
Again,
\[
1-\frac1p<1-\frac1{p+1},
\]
and since \(8>1\),
\[
8^{1-\frac1p}<8^{1-\frac1{p+1}}.
\]
Hence
\[
\rho_p(M)<\rho_{p+1}(M).
\]

Combining the two cases, for every integer \(n>6\) and every integer \(p>2\), there exists a matrix \(M\in \R^{n\times n}\) such that
\[
\boxed{\rho_p(M)<\rho_{p+1}(M)}.
\]
This proves Conjecture \(3\).
```